% Corps du contrôle 16 - Proportionnalité (partie 2 : pourcentages)

\bareme{
	\exbadge{1}{6} \quad
	\exbadge{2}{4} \quad
	\exbadge{3}{4} \quad
	\exbadge{4}{6} \quad
	\bonusbadge{2}
}

\begin{consignebox}
	\faIcon{info-circle} \textbf{CONSIGNES IMPORTANTES}
	\begin{itemize}[leftmargin=*,noitemsep,topsep=5pt]
		\item[\faIcon{check}] Justifiez les calculs (au moins une ligne de calcul par question).
		\item[\faIcon{check}] Donnez des unités (€, cm, km, g, \%).
		\item[\faIcon{check}] La calculatrice est interdite : choisissez des méthodes simples (10\%, 5\%, etc.).
	\end{itemize}
\end{consignebox}

\smallspace

\ExTemplate{1}{Pourcentages}{6}{calculator}{
\begin{enumerate}[label=\textcolor{primary}{\textbf{\alph*)}},leftmargin=*]
	\item Écrire sous forme de pourcentage : \(0{,}25\). \points{1}
	
	\anslines{1}{\(0{,}25 = 25\%\).}
	
	\item Calculer \(18\%\) de 50. \points{1}
	
	\anslines{2}{\(10\%\) de 50 = 5, donc \(18\% = 10\%+8\% = 5 + 4 = 9\).}
	
	\item Dans une classe de 28 élèves, \(25\%\) portent des lunettes. Combien d'élèves portent des lunettes ? \points{1}
	
	\anslines{2}{\(25\%\) = \(\dfrac{1}{4}\). Donc \(28 \div 4 = 7\) élèves.}
	
	\item Un article coûte 48 €. Il est en promotion : \(-15\%\). Calculer le montant de la réduction puis le prix payé. \points{3}
	
	\anslines{4}{
		\(10\%\) de 48 = 4,8 et \(5\%\) de 48 = 2,4 donc \(15\%\) de 48 = 7,2.
		Prix payé : \(48 - 7,2 = 40,8\) €.
	}
\end{enumerate}
}

\ExTemplate{2}{Tableau de proportionnalité}{4}{table}{
Un magasin vend des pommes. Le prix est proportionnel à la masse.

\smallspace

\[
\begin{array}{|c|c|c|c|c|}
\hline
\text{Masse (kg)} & 2 & 3 & 5 & 7\\
\hline
\text{Prix (€)} & 3{,}20 & 4{,}80 & \solution[2cm]{8} & \solution[2cm]{11{,}20}\\
\hline
\end{array}
\]

\begin{enumerate}[label=\textcolor{primary}{\textbf{\alph*)}},leftmargin=*]
	\item Quel est le prix de 1 kg de pommes ? \points{2}
	
	\anslines{2}{\(3{,}20 \div 2 = 1{,}60\) € pour 1 kg.}
	
	\item Expliquer comment trouver le prix de 7 kg. \points{2}
	
	\anslines{2}{On multiplie le prix de 1 kg par 7 : \(1{,}60\times 7 = 11{,}20\) €.}
\end{enumerate}
}

\ExTemplate{3}{Échelle}{4}{book}{
Une carte est à l'échelle \(1 : 25\,000\).

\begin{enumerate}[label=\textcolor{primary}{\textbf{\alph*)}},leftmargin=*]
	\item Sur la carte, la distance entre deux villes est de 6,4 cm. Quelle est la distance réelle (en km) ? \points{2}
	
	\anslines{3}{
		Distance réelle : \(6,4\times 25\,000 = 160\,000\) cm.
		Or \(100\,000\) cm = 1 km, donc \(160\,000\) cm = 1,6 km.
	}
	
	\item La distance réelle entre deux villages est de 3 km. Quelle distance cela représente-t-il sur la carte (en cm) ? \points{2}
	
	\anslines{2}{
		3 km = 300\,000 cm. Sur la carte : \(300\,000 \div 25\,000 = 12\) cm.
	}
\end{enumerate}
}

\begin{problembox}{\faIcon{utensils} Exercice 4 : Problème — Proportionnalité et pourcentages \points{6}}
	Pour faire \textbf{6 crêpes}, on utilise : \textbf{150 g} de farine, \textbf{2} œufs et \textbf{25 cL} de lait.
	
	\begin{enumerate}[label=\textcolor{primary}{\textbf{\alph*)}},leftmargin=*]
		\item Quantités pour faire \textbf{15 crêpes}. \points{4}
		
		\anslines{5}{
			On multiplie par \(15\div 6 = 2{,}5\).
			Farine : \(150\times 2{,}5 = 375\) g.
			Œufs : \(2\times 2{,}5 = 5\) œufs.
			Lait : \(25\times 2{,}5 = 62{,}5\) cL.
		}
		
		\item On décide d'ajouter \(+20\%\) de farine par rapport à la quantité calculée à la question a). Quelle masse de farine doit-on finalement utiliser ? \points{2}
		
		\anslines{3}{
			\(20\%\) de 375 = \(10\%\times 2\) = \(37{,}5\times 2 = 75\) g.
			Donc farine : \(375 + 75 = 450\) g.
		}
	\end{enumerate}
\end{problembox}

\exspace

\begin{exercisebox}{\faIcon{gift} Bonus \points{2}}
	Sur un graphique, une droite passe par l'origine et par le point \((4 ; 10)\).
	
	\smallspace
	
	\textbf{1.} Quel est le coefficient de proportionnalité ? \points{1}
	
	\anslines{1}{Coefficient : \(10 \div 4 = 2{,}5\).}
	
	\textbf{2.} Quelle est l'ordonnée du point de la droite d'abscisse 7 ? \points{1}
	
	\anslines{1}{\(y = 2{,}5\times 7 = 17{,}5\).}
\end{exercisebox}
